\chapter{May và Rủi}
\markboth{May và Rủi}{Luck and Misfortune}

\section{Một người lỡ chuyến xe và thoát khỏi một tai nạn lớn.}
\begin{truyen}{Một người lỡ chuyến xe và thoát khỏi một tai nạn lớn.}{Bad luck can hide a rescue.}
\chuhoa{A}{nh bực bội vì trễ giờ và nghĩ mình quá xui. Đến khi nghe tin chuyến xe ấy gặp sự cố, anh lặng đi. Có những điều rủi lúc đầu chỉ vì ta chưa biết hết câu chuyện.}
\textit{(He was annoyed at being late and felt deeply unlucky. When he later heard that the bus had an accident, he went silent. Some misfortunes seem bad only because we do not yet know the full story.)}
\end{truyen}
\begin{baihoc}
May rủi nhiều khi chỉ khác nhau ở chỗ ta biết câu chuyện đến đâu.
\textit{(Luck and misfortune often differ only in how much of the story we know.)}
\end{baihoc}
\begin{ghinhoanh}
\textbf{Short English to remember:}

Bad luck can hide a rescue.
\end{ghinhoanh}
\begin{tuhocnhanh}
1. Có chuyện nào lúc đầu em thấy xui nhưng về sau lại hóa may? \textit{(Has anything first felt unlucky and later turned out well?)}

2. Em có đang kết luận quá sớm về một chuyện không? \textit{(Are you judging any event too early?)}
\end{tuhocnhanh}
\ngancach

\section{Một người trúng thời nhưng không biết giữ mình.}
\begin{truyen}{Một người trúng thời nhưng không biết giữ mình.}{Luck without character can disappear quickly.}
\chuhoa{G}{ặp đúng thời điểm giúp anh đi lên nhanh. Nhưng vì nghĩ mình thật sự hơn người, anh chủ quan, coi thường người khác và lãng phí cơ hội. Vận may đưa anh lên, nhưng tính cách lại đẩy anh xuống.}
\textit{(The right timing helped him rise quickly. But believing himself superior, he became careless, dismissive, and wasteful. Luck lifted him up, yet character brought him down.)}
\end{truyen}
\begin{baihoc}
May mắn mở cửa, nhưng bản lĩnh mới giữ được cánh cửa ấy lâu.
\textit{(Luck opens doors, but character keeps them open.)}
\end{baihoc}
\begin{ghinhoanh}
\textbf{Short English to remember:}

Luck opens doors. Character keeps them open.
\end{ghinhoanh}
\begin{tuhocnhanh}
1. Em có đang nhầm may mắn với năng lực tuyệt đối không? \textit{(Are you mistaking good luck for total superiority?)}

2. Nếu gặp thời, em muốn giữ mình bằng điều gì? \textit{(If luck comes, what will help you stay grounded?)}
\end{tuhocnhanh}
\ngancach

\section{Một người rớt cơ hội rồi gặp con đường hợp hơn.}
\begin{truyen}{Một người rớt cơ hội rồi gặp con đường hợp hơn.}{Misfortune can redirect talent.}
\chuhoa{A}{nh tiếc đứt ruột vì hụt một suất mình từng rất mong. Nhưng chính quãng trống sau đó đưa anh đến công việc hợp khả năng và tính nết hơn nhiều. Điều ta gọi là xui đôi khi chỉ là đời đang bẻ lái.}
\textit{(He was deeply upset after missing a position he badly wanted. Yet the empty season afterward led him to work that fit his ability and personality much better. What we call bad luck is sometimes life changing direction.)}
\end{truyen}
\begin{baihoc}
Rủi ro trước mắt đôi khi là cách cuộc đời cứu ta khỏi một đường dài không hợp.
\textit{(A visible setback can sometimes save us from a long unsuitable road.)}
\end{baihoc}
\begin{ghinhoanh}
\textbf{Short English to remember:}

Misfortune can redirect talent.
\end{ghinhoanh}
\begin{tuhocnhanh}
1. Em có đang tiếc một cánh cửa đóng mà lẽ ra nên cảm ơn không? \textit{(Are you grieving a closed door you may one day thank?)}

2. Đời đang bẻ lái em về đâu? \textit{(Where might life be redirecting you?)}
\end{tuhocnhanh}
\ngancach

\section{Một người quá tin vận đỏ nên lười chuẩn bị.}
\begin{truyen}{Một người quá tin vận đỏ nên lười chuẩn bị.}{Luck should not replace effort.}
\chuhoa{V}{ài lần thuận lợi liên tiếp khiến anh nghĩ mình luôn gặp may. Vì thế anh chuẩn bị sơ sài, làm việc cẩu thả hơn. Khi kết quả xấu đến, anh mới hiểu may mắn không có nhiệm vụ bù cho sự cẩu thả của mình.}
\textit{(A few lucky breaks made him believe fortune would always follow. He began preparing carelessly and working loosely. When failure arrived, he learned that luck is not meant to cover our negligence.)}
\end{truyen}
\begin{baihoc}
May mắn là phần thêm vào, không phải nền móng để mình sống cẩu thả.
\textit{(Luck is an addition, not a foundation for careless living.)}
\end{baihoc}
\begin{ghinhoanh}
\textbf{Short English to remember:}

Luck should not replace effort.
\end{ghinhoanh}
\begin{tuhocnhanh}
1. Em có đang trông chờ vận may ở chỗ lẽ ra phải chuẩn bị không? \textit{(Are you waiting for luck where preparation is needed?)}

2. Điều gì em cần làm kỹ hơn ngay bây giờ? \textit{(What do you need to prepare more carefully now?)}
\end{tuhocnhanh}
\ngancach

\section{Một gia đình gặp biến cố nhưng nhờ đó xích lại gần nhau.}
\begin{truyen}{Một gia đình gặp biến cố nhưng nhờ đó xích lại gần nhau.}{Some hardship gathers people back together.}
\chuhoa{B}{ệnh tật và khó khăn làm cả nhà lo lắng. Nhưng cũng chính khoảng thời gian ấy khiến họ thôi bận bịu riêng, ngồi lại với nhau nhiều hơn, nói với nhau thật hơn. Có những điều không may làm hiện ra những điều quý bị lãng quên.}
\textit{(Illness and difficulty worried the whole family. But that season also pulled them out of separate busyness and made them sit together more, speak more honestly. Some hardships reveal precious things that had been forgotten.)}
\end{truyen}
\begin{baihoc}
Không phải mọi chuyện rủi đều chỉ lấy đi; có cái còn gọi những người xa nhau quay lại gần.
\textit{(Not every misfortune only takes away; some call distant people back together.)}
\end{baihoc}
\begin{ghinhoanh}
\textbf{Short English to remember:}

Some hardship gathers people back.
\end{ghinhoanh}
\begin{tuhocnhanh}
1. Khó khăn nào từng làm em thấy rõ điều quý hơn? \textit{(What hardship helped you see something more precious?)}

2. Điều quý đó hôm nay em còn giữ không? \textit{(Are you still keeping that precious thing today?)}
\end{tuhocnhanh}
\ngancach